\section{Latar Belakang}

Pariwisata merupakan salah satu sektor yang berperan penting dalam meningkatkan pertumbuhan ekonomi daerah serta mendukung kesejahteraan masyarakat. Selain memberikan kontribusi terhadap pendapatan daerah, sektor pariwisata juga mampu mendorong perkembangan infrastruktur dan membuka peluang usaha bagi masyarakat sekitar \cite{afendi2023}.

Perkembangan teknologi informasi telah mengubah cara wisatawan dalam mencari informasi dan memberikan penilaian terhadap suatu destinasi wisata. Saat ini, platform digital seperti Google Reviews menjadi salah satu media yang banyak digunakan wisatawan untuk menyampaikan pengalaman, kritik, dan saran setelah berkunjung ke suatu objek wisata. Informasi yang terkandung dalam ulasan tersebut dapat menjadi sumber masukan yang berharga bagi pengelola destinasi wisata \cite{ramadhani2024}.

Kabupaten Bangkalan memiliki berbagai destinasi wisata alam dan religi yang menarik minat wisatawan. Namun, banyaknya ulasan yang terus bertambah menyebabkan proses analisis secara manual menjadi kurang efektif. Oleh karena itu, diperlukan suatu pendekatan yang mampu membantu pengelola memahami persepsi pengunjung secara lebih cepat dan sistematis sehingga dapat digunakan sebagai dasar dalam meningkatkan kualitas layanan dan pengembangan destinasi wisata.